\chapter{Anexos}
\label{cap:anexos}

Este capitulo consolida los artefactos tecnicos que soportan la trazabilidad, reproducibilidad y auditoria del estudio. Los anexos complementan los capitulos principales y permiten verificar configuracion, ejecucion, procesamiento de datos y resultados estadisticos.

\section{A. Configuracion experimental completa}
\label{anx:config}

\subsection{A.1 Entorno de ejecucion}
\begin{itemize}
  \item Hardware anfitrion: AMD Ryzen 5 3600, 16 GB RAM.
  \item Sistema anfitrion: Windows con WSL2.
  \item VM/WSL2 asignado al experimento: 4 vCPU, 12 GB RAM, 2 GB swap.
  \item Sistema base: Ubuntu 20.04.
\end{itemize}

\subsection{A.2 Stack de herramientas}
\begin{itemize}
  \item Kubernetes: MicroK8s.
  \item Pruebas de carga: k6 v0.50.0.
  \item Observabilidad: Prometheus + Grafana.
  \item Orquestacion de corridas: scripts en Bash/Python del repositorio.
\end{itemize}

\subsection{A.3 Parametros operativos}
\begin{itemize}
  \item Niveles de carga: VUS = 1, 5, 10, 20.
  \item Controles evaluados: C1, C2, C3, C4.
  \item Variantes por control: 3.
  \item Replicas por celda: 8.
  \item Bloqueo temporal: 2 dias de ejecucion.
\end{itemize}

\section{B. Matriz de diseno experimental}
\label{anx:matriz-diseno}

\subsection{B.1 Estructura factorial}
Dise\~no factorial completo: $4 \times 3 \times 4 = 48$ celdas experimentales.

\subsection{B.2 Definicion de factores y niveles}
\begin{itemize}
  \item Factor 1 (Control): C1 API Gateway, C2 mTLS, C3 NetworkPolicy, C4 Rate Limiting.
  \item Factor 2 (Variante): 3 niveles por cada control.
  \item Factor 3 (Carga): VUS 1, 5, 10, 20.
\end{itemize}

\subsection{B.3 Matriz codificada de corridas}
Incluir tabla completa con columnas minimas:
\begin{itemize}
  \item id\_run
  \item day\_block
  \item control
  \item variant
  \item vus
  \item replica
  \item seed (si aplica)
\end{itemize}

\section{C. Tablas completas de resultados}
\label{anx:tablas-resultados}

Incluir resultados agregados y por replica para cada control y variante. Columnas recomendadas:
\begin{itemize}
  \item avg\_ms
  \item p95\_ms
  \item err\_pct
  \item rps
  \item cpu\_mcores
  \item mem\_mib
\end{itemize}

\subsection{C.1 Resultados por control}
\begin{itemize}
  \item C1: API Gateway.
  \item C2: mTLS.
  \item C3: NetworkPolicy.
  \item C4: Rate Limiting.
\end{itemize}

\section{D. Salidas ANOVA completas}
\label{anx:anova}

\subsection{D.1 Especificacion del modelo}
\begin{equation*}
y \sim C(\text{control}) * C(\text{variant\_level}) * C(\text{vus}) + C(\text{day})
\end{equation*}

\subsection{D.2 Tablas ANOVA por metrica}
Para cada metrica (avg\_ms, p95\_ms, err\_pct, rps, cpu\_mcores, mem\_mib) incluir:
\begin{itemize}
  \item Source
  \item df
  \item sum\_sq
  \item mean\_sq
  \item F
  \item p
  \item eta\_p2
\end{itemize}

\subsection{D.3 Archivos fuente sugeridos}
\begin{itemize}
  \item Testing/results/scaling\_tests/anova\_factorial\_full\_blocked\_by\_day\_20260510.csv
  \item Testing/results/scaling\_tests/anova\_factorial\_matrix\_20260510.csv
\end{itemize}

\section{E. Validacion estadistica y calidad de datos}
\label{anx:validacion}

\subsection{E.1 Verificaciones de consistencia}
\begin{itemize}
  \item Cobertura total de celdas: 48/48.
  \item Replicas por celda: 8.
  \item Revisiones de valores atipicos y corridas fallidas.
\end{itemize}

\subsection{E.2 Indicadores de estabilidad}
\begin{itemize}
  \item Coeficiente de variacion por metrica y por celda.
  \item Trazabilidad de corridas repetidas entre dias.
\end{itemize}

\section{F. Protocolo de reproducibilidad}
\label{anx:protocolo}

\subsection{F.1 Flujo de ejecucion}
\begin{enumerate}
  \item Preparacion del entorno y validacion de dependencias.
  \item Despliegue de la variante a evaluar.
  \item Warmup.
  \item Ejecucion de carga k6.
  \item Recoleccion de metricas de sistema.
  \item Cooldown y limpieza.
  \item Registro de artefactos de salida.
\end{enumerate}

\subsection{F.2 Comandos y scripts}
Incluir comandos exactos y scripts de referencia usados en campana.

\section{G. Artefactos de despliegue y configuracion}
\label{anx:artefactos}

Incluir manifiestos y configuraciones:
\begin{itemize}
  \item YAML de Gateway/Ingress.
  \item Politicas de mTLS.
  \item NetworkPolicy (baseline, basic, strict).
  \item Configuracion de Rate Limiting (baseline, moderate, strict).
\end{itemize}

\section{H. Diccionario de variables}
\label{anx:diccionario}

\begin{itemize}
  \item avg\_ms: latencia promedio por corrida (ms).
  \item p95\_ms: percentil 95 de latencia por corrida (ms).
  \item err\_pct: porcentaje de respuestas con error.
  \item rps: solicitudes por segundo.
  \item cpu\_mcores: uso promedio de CPU (millicores).
  \item mem\_mib: uso promedio de memoria (MiB).
  \item login\_success\_total: total de logins validos durante la corrida.
\end{itemize}

\section{I. Evidencia visual de operacion}
\label{anx:evidencia-visual}

Agregar capturas de paneles de Grafana y, cuando aplique, snapshots de estado del cluster durante corridas representativas:
\begin{itemize}
  \item Paneles de latencia y throughput.
  \item Paneles de CPU y memoria.
  \item Paneles de error rate.
  \item Evidencia de preservacion de usuarios logueados.
\end{itemize}

\section{Nota de uso}

Si el documento principal usa archivos separados por capitulo, este archivo puede integrarse con:
\begin{verbatim}
\chapter{Anexos}
\label{cap:anexos}

Este capitulo consolida los artefactos tecnicos que soportan la trazabilidad, reproducibilidad y auditoria del estudio. Los anexos complementan los capitulos principales y permiten verificar configuracion, ejecucion, procesamiento de datos y resultados estadisticos.

\section{A. Configuracion experimental completa}
\label{anx:config}

\subsection{A.1 Entorno de ejecucion}
\begin{itemize}
  \item Hardware anfitrion: AMD Ryzen 5 3600, 16 GB RAM.
  \item Sistema anfitrion: Windows con WSL2.
  \item VM/WSL2 asignado al experimento: 4 vCPU, 12 GB RAM, 2 GB swap.
  \item Sistema base: Ubuntu 20.04.
\end{itemize}

\subsection{A.2 Stack de herramientas}
\begin{itemize}
  \item Kubernetes: MicroK8s.
  \item Pruebas de carga: k6 v0.50.0.
  \item Observabilidad: Prometheus + Grafana.
  \item Orquestacion de corridas: scripts en Bash/Python del repositorio.
\end{itemize}

\subsection{A.3 Parametros operativos}
\begin{itemize}
  \item Niveles de carga: VUS = 1, 5, 10, 20.
  \item Controles evaluados: C1, C2, C3, C4.
  \item Variantes por control: 3.
  \item Replicas por celda: 8.
  \item Bloqueo temporal: 2 dias de ejecucion.
\end{itemize}

\section{B. Matriz de diseno experimental}
\label{anx:matriz-diseno}

\subsection{B.1 Estructura factorial}
Dise\~no factorial completo: $4 \times 3 \times 4 = 48$ celdas experimentales.

\subsection{B.2 Definicion de factores y niveles}
\begin{itemize}
  \item Factor 1 (Control): C1 API Gateway, C2 mTLS, C3 NetworkPolicy, C4 Rate Limiting.
  \item Factor 2 (Variante): 3 niveles por cada control.
  \item Factor 3 (Carga): VUS 1, 5, 10, 20.
\end{itemize}

\subsection{B.3 Matriz codificada de corridas}
Incluir tabla completa con columnas minimas:
\begin{itemize}
  \item id\_run
  \item day\_block
  \item control
  \item variant
  \item vus
  \item replica
  \item seed (si aplica)
\end{itemize}

\section{C. Tablas completas de resultados}
\label{anx:tablas-resultados}

Incluir resultados agregados y por replica para cada control y variante. Columnas recomendadas:
\begin{itemize}
  \item avg\_ms
  \item p95\_ms
  \item err\_pct
  \item rps
  \item cpu\_mcores
  \item mem\_mib
\end{itemize}

\subsection{C.1 Resultados por control}
\begin{itemize}
  \item C1: API Gateway.
  \item C2: mTLS.
  \item C3: NetworkPolicy.
  \item C4: Rate Limiting.
\end{itemize}

\section{D. Salidas ANOVA completas}
\label{anx:anova}

\subsection{D.1 Especificacion del modelo}
\begin{equation*}
y \sim C(\text{control}) * C(\text{variant\_level}) * C(\text{vus}) + C(\text{day})
\end{equation*}

\subsection{D.2 Tablas ANOVA por metrica}
Para cada metrica (avg\_ms, p95\_ms, err\_pct, rps, cpu\_mcores, mem\_mib) incluir:
\begin{itemize}
  \item Source
  \item df
  \item sum\_sq
  \item mean\_sq
  \item F
  \item p
  \item eta\_p2
\end{itemize}

\subsection{D.3 Archivos fuente sugeridos}
\begin{itemize}
  \item Testing/results/scaling\_tests/anova\_factorial\_full\_blocked\_by\_day\_20260510.csv
  \item Testing/results/scaling\_tests/anova\_factorial\_matrix\_20260510.csv
\end{itemize}

\section{E. Validacion estadistica y calidad de datos}
\label{anx:validacion}

\subsection{E.1 Verificaciones de consistencia}
\begin{itemize}
  \item Cobertura total de celdas: 48/48.
  \item Replicas por celda: 8.
  \item Revisiones de valores atipicos y corridas fallidas.
\end{itemize}

\subsection{E.2 Indicadores de estabilidad}
\begin{itemize}
  \item Coeficiente de variacion por metrica y por celda.
  \item Trazabilidad de corridas repetidas entre dias.
\end{itemize}

\section{F. Protocolo de reproducibilidad}
\label{anx:protocolo}

\subsection{F.1 Flujo de ejecucion}
\begin{enumerate}
  \item Preparacion del entorno y validacion de dependencias.
  \item Despliegue de la variante a evaluar.
  \item Warmup.
  \item Ejecucion de carga k6.
  \item Recoleccion de metricas de sistema.
  \item Cooldown y limpieza.
  \item Registro de artefactos de salida.
\end{enumerate}

\subsection{F.2 Comandos y scripts}
Incluir comandos exactos y scripts de referencia usados en campana.

\section{G. Artefactos de despliegue y configuracion}
\label{anx:artefactos}

Incluir manifiestos y configuraciones:
\begin{itemize}
  \item YAML de Gateway/Ingress.
  \item Politicas de mTLS.
  \item NetworkPolicy (baseline, basic, strict).
  \item Configuracion de Rate Limiting (baseline, moderate, strict).
\end{itemize}

\section{H. Diccionario de variables}
\label{anx:diccionario}

\begin{itemize}
  \item avg\_ms: latencia promedio por corrida (ms).
  \item p95\_ms: percentil 95 de latencia por corrida (ms).
  \item err\_pct: porcentaje de respuestas con error.
  \item rps: solicitudes por segundo.
  \item cpu\_mcores: uso promedio de CPU (millicores).
  \item mem\_mib: uso promedio de memoria (MiB).
  \item login\_success\_total: total de logins validos durante la corrida.
\end{itemize}

\section{I. Evidencia visual de operacion}
\label{anx:evidencia-visual}

Agregar capturas de paneles de Grafana y, cuando aplique, snapshots de estado del cluster durante corridas representativas:
\begin{itemize}
  \item Paneles de latencia y throughput.
  \item Paneles de CPU y memoria.
  \item Paneles de error rate.
  \item Evidencia de preservacion de usuarios logueados.
\end{itemize}

\section{Nota de uso}

Si el documento principal usa archivos separados por capitulo, este archivo puede integrarse con:
\begin{verbatim}
\chapter{Anexos}
\label{cap:anexos}

Este capitulo consolida los artefactos tecnicos que soportan la trazabilidad, reproducibilidad y auditoria del estudio. Los anexos complementan los capitulos principales y permiten verificar configuracion, ejecucion, procesamiento de datos y resultados estadisticos.

\section{A. Configuracion experimental completa}
\label{anx:config}

\subsection{A.1 Entorno de ejecucion}
\begin{itemize}
  \item Hardware anfitrion: AMD Ryzen 5 3600, 16 GB RAM.
  \item Sistema anfitrion: Windows con WSL2.
  \item VM/WSL2 asignado al experimento: 4 vCPU, 12 GB RAM, 2 GB swap.
  \item Sistema base: Ubuntu 20.04.
\end{itemize}

\subsection{A.2 Stack de herramientas}
\begin{itemize}
  \item Kubernetes: MicroK8s.
  \item Pruebas de carga: k6 v0.50.0.
  \item Observabilidad: Prometheus + Grafana.
  \item Orquestacion de corridas: scripts en Bash/Python del repositorio.
\end{itemize}

\subsection{A.3 Parametros operativos}
\begin{itemize}
  \item Niveles de carga: VUS = 1, 5, 10, 20.
  \item Controles evaluados: C1, C2, C3, C4.
  \item Variantes por control: 3.
  \item Replicas por celda: 8.
  \item Bloqueo temporal: 2 dias de ejecucion.
\end{itemize}

\section{B. Matriz de diseno experimental}
\label{anx:matriz-diseno}

\subsection{B.1 Estructura factorial}
Dise\~no factorial completo: $4 \times 3 \times 4 = 48$ celdas experimentales.

\subsection{B.2 Definicion de factores y niveles}
\begin{itemize}
  \item Factor 1 (Control): C1 API Gateway, C2 mTLS, C3 NetworkPolicy, C4 Rate Limiting.
  \item Factor 2 (Variante): 3 niveles por cada control.
  \item Factor 3 (Carga): VUS 1, 5, 10, 20.
\end{itemize}

\subsection{B.3 Matriz codificada de corridas}
Incluir tabla completa con columnas minimas:
\begin{itemize}
  \item id\_run
  \item day\_block
  \item control
  \item variant
  \item vus
  \item replica
  \item seed (si aplica)
\end{itemize}

\section{C. Tablas completas de resultados}
\label{anx:tablas-resultados}

Incluir resultados agregados y por replica para cada control y variante. Columnas recomendadas:
\begin{itemize}
  \item avg\_ms
  \item p95\_ms
  \item err\_pct
  \item rps
  \item cpu\_mcores
  \item mem\_mib
\end{itemize}

\subsection{C.1 Resultados por control}
\begin{itemize}
  \item C1: API Gateway.
  \item C2: mTLS.
  \item C3: NetworkPolicy.
  \item C4: Rate Limiting.
\end{itemize}

\section{D. Salidas ANOVA completas}
\label{anx:anova}

\subsection{D.1 Especificacion del modelo}
\begin{equation*}
y \sim C(\text{control}) * C(\text{variant\_level}) * C(\text{vus}) + C(\text{day})
\end{equation*}

\subsection{D.2 Tablas ANOVA por metrica}
Para cada metrica (avg\_ms, p95\_ms, err\_pct, rps, cpu\_mcores, mem\_mib) incluir:
\begin{itemize}
  \item Source
  \item df
  \item sum\_sq
  \item mean\_sq
  \item F
  \item p
  \item eta\_p2
\end{itemize}

\subsection{D.3 Archivos fuente sugeridos}
\begin{itemize}
  \item Testing/results/scaling\_tests/anova\_factorial\_full\_blocked\_by\_day\_20260510.csv
  \item Testing/results/scaling\_tests/anova\_factorial\_matrix\_20260510.csv
\end{itemize}

\section{E. Validacion estadistica y calidad de datos}
\label{anx:validacion}

\subsection{E.1 Verificaciones de consistencia}
\begin{itemize}
  \item Cobertura total de celdas: 48/48.
  \item Replicas por celda: 8.
  \item Revisiones de valores atipicos y corridas fallidas.
\end{itemize}

\subsection{E.2 Indicadores de estabilidad}
\begin{itemize}
  \item Coeficiente de variacion por metrica y por celda.
  \item Trazabilidad de corridas repetidas entre dias.
\end{itemize}

\section{F. Protocolo de reproducibilidad}
\label{anx:protocolo}

\subsection{F.1 Flujo de ejecucion}
\begin{enumerate}
  \item Preparacion del entorno y validacion de dependencias.
  \item Despliegue de la variante a evaluar.
  \item Warmup.
  \item Ejecucion de carga k6.
  \item Recoleccion de metricas de sistema.
  \item Cooldown y limpieza.
  \item Registro de artefactos de salida.
\end{enumerate}

\subsection{F.2 Comandos y scripts}
Incluir comandos exactos y scripts de referencia usados en campana.

\section{G. Artefactos de despliegue y configuracion}
\label{anx:artefactos}

Incluir manifiestos y configuraciones:
\begin{itemize}
  \item YAML de Gateway/Ingress.
  \item Politicas de mTLS.
  \item NetworkPolicy (baseline, basic, strict).
  \item Configuracion de Rate Limiting (baseline, moderate, strict).
\end{itemize}

\section{H. Diccionario de variables}
\label{anx:diccionario}

\begin{itemize}
  \item avg\_ms: latencia promedio por corrida (ms).
  \item p95\_ms: percentil 95 de latencia por corrida (ms).
  \item err\_pct: porcentaje de respuestas con error.
  \item rps: solicitudes por segundo.
  \item cpu\_mcores: uso promedio de CPU (millicores).
  \item mem\_mib: uso promedio de memoria (MiB).
  \item login\_success\_total: total de logins validos durante la corrida.
\end{itemize}

\section{I. Evidencia visual de operacion}
\label{anx:evidencia-visual}

Agregar capturas de paneles de Grafana y, cuando aplique, snapshots de estado del cluster durante corridas representativas:
\begin{itemize}
  \item Paneles de latencia y throughput.
  \item Paneles de CPU y memoria.
  \item Paneles de error rate.
  \item Evidencia de preservacion de usuarios logueados.
\end{itemize}

\section{Nota de uso}

Si el documento principal usa archivos separados por capitulo, este archivo puede integrarse con:
\begin{verbatim}
\chapter{Anexos}
\label{cap:anexos}

Este capitulo consolida los artefactos tecnicos que soportan la trazabilidad, reproducibilidad y auditoria del estudio. Los anexos complementan los capitulos principales y permiten verificar configuracion, ejecucion, procesamiento de datos y resultados estadisticos.

\section{A. Configuracion experimental completa}
\label{anx:config}

\subsection{A.1 Entorno de ejecucion}
\begin{itemize}
  \item Hardware anfitrion: AMD Ryzen 5 3600, 16 GB RAM.
  \item Sistema anfitrion: Windows con WSL2.
  \item VM/WSL2 asignado al experimento: 4 vCPU, 12 GB RAM, 2 GB swap.
  \item Sistema base: Ubuntu 20.04.
\end{itemize}

\subsection{A.2 Stack de herramientas}
\begin{itemize}
  \item Kubernetes: MicroK8s.
  \item Pruebas de carga: k6 v0.50.0.
  \item Observabilidad: Prometheus + Grafana.
  \item Orquestacion de corridas: scripts en Bash/Python del repositorio.
\end{itemize}

\subsection{A.3 Parametros operativos}
\begin{itemize}
  \item Niveles de carga: VUS = 1, 5, 10, 20.
  \item Controles evaluados: C1, C2, C3, C4.
  \item Variantes por control: 3.
  \item Replicas por celda: 8.
  \item Bloqueo temporal: 2 dias de ejecucion.
\end{itemize}

\section{B. Matriz de diseno experimental}
\label{anx:matriz-diseno}

\subsection{B.1 Estructura factorial}
Dise\~no factorial completo: $4 \times 3 \times 4 = 48$ celdas experimentales.

\subsection{B.2 Definicion de factores y niveles}
\begin{itemize}
  \item Factor 1 (Control): C1 API Gateway, C2 mTLS, C3 NetworkPolicy, C4 Rate Limiting.
  \item Factor 2 (Variante): 3 niveles por cada control.
  \item Factor 3 (Carga): VUS 1, 5, 10, 20.
\end{itemize}

\subsection{B.3 Matriz codificada de corridas}
Incluir tabla completa con columnas minimas:
\begin{itemize}
  \item id\_run
  \item day\_block
  \item control
  \item variant
  \item vus
  \item replica
  \item seed (si aplica)
\end{itemize}

\section{C. Tablas completas de resultados}
\label{anx:tablas-resultados}

Incluir resultados agregados y por replica para cada control y variante. Columnas recomendadas:
\begin{itemize}
  \item avg\_ms
  \item p95\_ms
  \item err\_pct
  \item rps
  \item cpu\_mcores
  \item mem\_mib
\end{itemize}

\subsection{C.1 Resultados por control}
\begin{itemize}
  \item C1: API Gateway.
  \item C2: mTLS.
  \item C3: NetworkPolicy.
  \item C4: Rate Limiting.
\end{itemize}

\section{D. Salidas ANOVA completas}
\label{anx:anova}

\subsection{D.1 Especificacion del modelo}
\begin{equation*}
y \sim C(\text{control}) * C(\text{variant\_level}) * C(\text{vus}) + C(\text{day})
\end{equation*}

\subsection{D.2 Tablas ANOVA por metrica}
Para cada metrica (avg\_ms, p95\_ms, err\_pct, rps, cpu\_mcores, mem\_mib) incluir:
\begin{itemize}
  \item Source
  \item df
  \item sum\_sq
  \item mean\_sq
  \item F
  \item p
  \item eta\_p2
\end{itemize}

\subsection{D.3 Archivos fuente sugeridos}
\begin{itemize}
  \item Testing/results/scaling\_tests/anova\_factorial\_full\_blocked\_by\_day\_20260510.csv
  \item Testing/results/scaling\_tests/anova\_factorial\_matrix\_20260510.csv
\end{itemize}

\section{E. Validacion estadistica y calidad de datos}
\label{anx:validacion}

\subsection{E.1 Verificaciones de consistencia}
\begin{itemize}
  \item Cobertura total de celdas: 48/48.
  \item Replicas por celda: 8.
  \item Revisiones de valores atipicos y corridas fallidas.
\end{itemize}

\subsection{E.2 Indicadores de estabilidad}
\begin{itemize}
  \item Coeficiente de variacion por metrica y por celda.
  \item Trazabilidad de corridas repetidas entre dias.
\end{itemize}

\section{F. Protocolo de reproducibilidad}
\label{anx:protocolo}

\subsection{F.1 Flujo de ejecucion}
\begin{enumerate}
  \item Preparacion del entorno y validacion de dependencias.
  \item Despliegue de la variante a evaluar.
  \item Warmup.
  \item Ejecucion de carga k6.
  \item Recoleccion de metricas de sistema.
  \item Cooldown y limpieza.
  \item Registro de artefactos de salida.
\end{enumerate}

\subsection{F.2 Comandos y scripts}
Incluir comandos exactos y scripts de referencia usados en campana.

\section{G. Artefactos de despliegue y configuracion}
\label{anx:artefactos}

Incluir manifiestos y configuraciones:
\begin{itemize}
  \item YAML de Gateway/Ingress.
  \item Politicas de mTLS.
  \item NetworkPolicy (baseline, basic, strict).
  \item Configuracion de Rate Limiting (baseline, moderate, strict).
\end{itemize}

\section{H. Diccionario de variables}
\label{anx:diccionario}

\begin{itemize}
  \item avg\_ms: latencia promedio por corrida (ms).
  \item p95\_ms: percentil 95 de latencia por corrida (ms).
  \item err\_pct: porcentaje de respuestas con error.
  \item rps: solicitudes por segundo.
  \item cpu\_mcores: uso promedio de CPU (millicores).
  \item mem\_mib: uso promedio de memoria (MiB).
  \item login\_success\_total: total de logins validos durante la corrida.
\end{itemize}

\section{I. Evidencia visual de operacion}
\label{anx:evidencia-visual}

Agregar capturas de paneles de Grafana y, cuando aplique, snapshots de estado del cluster durante corridas representativas:
\begin{itemize}
  \item Paneles de latencia y throughput.
  \item Paneles de CPU y memoria.
  \item Paneles de error rate.
  \item Evidencia de preservacion de usuarios logueados.
\end{itemize}

\section{Nota de uso}

Si el documento principal usa archivos separados por capitulo, este archivo puede integrarse con:
\begin{verbatim}
\input{documentecionFinal/anexos}
\end{verbatim}

o bien, si el archivo principal esta en la misma carpeta:
\begin{verbatim}
\input{anexos}
\end{verbatim}

\end{verbatim}

o bien, si el archivo principal esta en la misma carpeta:
\begin{verbatim}
\chapter{Anexos}
\label{cap:anexos}

Este capitulo consolida los artefactos tecnicos que soportan la trazabilidad, reproducibilidad y auditoria del estudio. Los anexos complementan los capitulos principales y permiten verificar configuracion, ejecucion, procesamiento de datos y resultados estadisticos.

\section{A. Configuracion experimental completa}
\label{anx:config}

\subsection{A.1 Entorno de ejecucion}
\begin{itemize}
  \item Hardware anfitrion: AMD Ryzen 5 3600, 16 GB RAM.
  \item Sistema anfitrion: Windows con WSL2.
  \item VM/WSL2 asignado al experimento: 4 vCPU, 12 GB RAM, 2 GB swap.
  \item Sistema base: Ubuntu 20.04.
\end{itemize}

\subsection{A.2 Stack de herramientas}
\begin{itemize}
  \item Kubernetes: MicroK8s.
  \item Pruebas de carga: k6 v0.50.0.
  \item Observabilidad: Prometheus + Grafana.
  \item Orquestacion de corridas: scripts en Bash/Python del repositorio.
\end{itemize}

\subsection{A.3 Parametros operativos}
\begin{itemize}
  \item Niveles de carga: VUS = 1, 5, 10, 20.
  \item Controles evaluados: C1, C2, C3, C4.
  \item Variantes por control: 3.
  \item Replicas por celda: 8.
  \item Bloqueo temporal: 2 dias de ejecucion.
\end{itemize}

\section{B. Matriz de diseno experimental}
\label{anx:matriz-diseno}

\subsection{B.1 Estructura factorial}
Dise\~no factorial completo: $4 \times 3 \times 4 = 48$ celdas experimentales.

\subsection{B.2 Definicion de factores y niveles}
\begin{itemize}
  \item Factor 1 (Control): C1 API Gateway, C2 mTLS, C3 NetworkPolicy, C4 Rate Limiting.
  \item Factor 2 (Variante): 3 niveles por cada control.
  \item Factor 3 (Carga): VUS 1, 5, 10, 20.
\end{itemize}

\subsection{B.3 Matriz codificada de corridas}
Incluir tabla completa con columnas minimas:
\begin{itemize}
  \item id\_run
  \item day\_block
  \item control
  \item variant
  \item vus
  \item replica
  \item seed (si aplica)
\end{itemize}

\section{C. Tablas completas de resultados}
\label{anx:tablas-resultados}

Incluir resultados agregados y por replica para cada control y variante. Columnas recomendadas:
\begin{itemize}
  \item avg\_ms
  \item p95\_ms
  \item err\_pct
  \item rps
  \item cpu\_mcores
  \item mem\_mib
\end{itemize}

\subsection{C.1 Resultados por control}
\begin{itemize}
  \item C1: API Gateway.
  \item C2: mTLS.
  \item C3: NetworkPolicy.
  \item C4: Rate Limiting.
\end{itemize}

\section{D. Salidas ANOVA completas}
\label{anx:anova}

\subsection{D.1 Especificacion del modelo}
\begin{equation*}
y \sim C(\text{control}) * C(\text{variant\_level}) * C(\text{vus}) + C(\text{day})
\end{equation*}

\subsection{D.2 Tablas ANOVA por metrica}
Para cada metrica (avg\_ms, p95\_ms, err\_pct, rps, cpu\_mcores, mem\_mib) incluir:
\begin{itemize}
  \item Source
  \item df
  \item sum\_sq
  \item mean\_sq
  \item F
  \item p
  \item eta\_p2
\end{itemize}

\subsection{D.3 Archivos fuente sugeridos}
\begin{itemize}
  \item Testing/results/scaling\_tests/anova\_factorial\_full\_blocked\_by\_day\_20260510.csv
  \item Testing/results/scaling\_tests/anova\_factorial\_matrix\_20260510.csv
\end{itemize}

\section{E. Validacion estadistica y calidad de datos}
\label{anx:validacion}

\subsection{E.1 Verificaciones de consistencia}
\begin{itemize}
  \item Cobertura total de celdas: 48/48.
  \item Replicas por celda: 8.
  \item Revisiones de valores atipicos y corridas fallidas.
\end{itemize}

\subsection{E.2 Indicadores de estabilidad}
\begin{itemize}
  \item Coeficiente de variacion por metrica y por celda.
  \item Trazabilidad de corridas repetidas entre dias.
\end{itemize}

\section{F. Protocolo de reproducibilidad}
\label{anx:protocolo}

\subsection{F.1 Flujo de ejecucion}
\begin{enumerate}
  \item Preparacion del entorno y validacion de dependencias.
  \item Despliegue de la variante a evaluar.
  \item Warmup.
  \item Ejecucion de carga k6.
  \item Recoleccion de metricas de sistema.
  \item Cooldown y limpieza.
  \item Registro de artefactos de salida.
\end{enumerate}

\subsection{F.2 Comandos y scripts}
Incluir comandos exactos y scripts de referencia usados en campana.

\section{G. Artefactos de despliegue y configuracion}
\label{anx:artefactos}

Incluir manifiestos y configuraciones:
\begin{itemize}
  \item YAML de Gateway/Ingress.
  \item Politicas de mTLS.
  \item NetworkPolicy (baseline, basic, strict).
  \item Configuracion de Rate Limiting (baseline, moderate, strict).
\end{itemize}

\section{H. Diccionario de variables}
\label{anx:diccionario}

\begin{itemize}
  \item avg\_ms: latencia promedio por corrida (ms).
  \item p95\_ms: percentil 95 de latencia por corrida (ms).
  \item err\_pct: porcentaje de respuestas con error.
  \item rps: solicitudes por segundo.
  \item cpu\_mcores: uso promedio de CPU (millicores).
  \item mem\_mib: uso promedio de memoria (MiB).
  \item login\_success\_total: total de logins validos durante la corrida.
\end{itemize}

\section{I. Evidencia visual de operacion}
\label{anx:evidencia-visual}

Agregar capturas de paneles de Grafana y, cuando aplique, snapshots de estado del cluster durante corridas representativas:
\begin{itemize}
  \item Paneles de latencia y throughput.
  \item Paneles de CPU y memoria.
  \item Paneles de error rate.
  \item Evidencia de preservacion de usuarios logueados.
\end{itemize}

\section{Nota de uso}

Si el documento principal usa archivos separados por capitulo, este archivo puede integrarse con:
\begin{verbatim}
\input{documentecionFinal/anexos}
\end{verbatim}

o bien, si el archivo principal esta en la misma carpeta:
\begin{verbatim}
\input{anexos}
\end{verbatim}

\end{verbatim}

\end{verbatim}

o bien, si el archivo principal esta en la misma carpeta:
\begin{verbatim}
\chapter{Anexos}
\label{cap:anexos}

Este capitulo consolida los artefactos tecnicos que soportan la trazabilidad, reproducibilidad y auditoria del estudio. Los anexos complementan los capitulos principales y permiten verificar configuracion, ejecucion, procesamiento de datos y resultados estadisticos.

\section{A. Configuracion experimental completa}
\label{anx:config}

\subsection{A.1 Entorno de ejecucion}
\begin{itemize}
  \item Hardware anfitrion: AMD Ryzen 5 3600, 16 GB RAM.
  \item Sistema anfitrion: Windows con WSL2.
  \item VM/WSL2 asignado al experimento: 4 vCPU, 12 GB RAM, 2 GB swap.
  \item Sistema base: Ubuntu 20.04.
\end{itemize}

\subsection{A.2 Stack de herramientas}
\begin{itemize}
  \item Kubernetes: MicroK8s.
  \item Pruebas de carga: k6 v0.50.0.
  \item Observabilidad: Prometheus + Grafana.
  \item Orquestacion de corridas: scripts en Bash/Python del repositorio.
\end{itemize}

\subsection{A.3 Parametros operativos}
\begin{itemize}
  \item Niveles de carga: VUS = 1, 5, 10, 20.
  \item Controles evaluados: C1, C2, C3, C4.
  \item Variantes por control: 3.
  \item Replicas por celda: 8.
  \item Bloqueo temporal: 2 dias de ejecucion.
\end{itemize}

\section{B. Matriz de diseno experimental}
\label{anx:matriz-diseno}

\subsection{B.1 Estructura factorial}
Dise\~no factorial completo: $4 \times 3 \times 4 = 48$ celdas experimentales.

\subsection{B.2 Definicion de factores y niveles}
\begin{itemize}
  \item Factor 1 (Control): C1 API Gateway, C2 mTLS, C3 NetworkPolicy, C4 Rate Limiting.
  \item Factor 2 (Variante): 3 niveles por cada control.
  \item Factor 3 (Carga): VUS 1, 5, 10, 20.
\end{itemize}

\subsection{B.3 Matriz codificada de corridas}
Incluir tabla completa con columnas minimas:
\begin{itemize}
  \item id\_run
  \item day\_block
  \item control
  \item variant
  \item vus
  \item replica
  \item seed (si aplica)
\end{itemize}

\section{C. Tablas completas de resultados}
\label{anx:tablas-resultados}

Incluir resultados agregados y por replica para cada control y variante. Columnas recomendadas:
\begin{itemize}
  \item avg\_ms
  \item p95\_ms
  \item err\_pct
  \item rps
  \item cpu\_mcores
  \item mem\_mib
\end{itemize}

\subsection{C.1 Resultados por control}
\begin{itemize}
  \item C1: API Gateway.
  \item C2: mTLS.
  \item C3: NetworkPolicy.
  \item C4: Rate Limiting.
\end{itemize}

\section{D. Salidas ANOVA completas}
\label{anx:anova}

\subsection{D.1 Especificacion del modelo}
\begin{equation*}
y \sim C(\text{control}) * C(\text{variant\_level}) * C(\text{vus}) + C(\text{day})
\end{equation*}

\subsection{D.2 Tablas ANOVA por metrica}
Para cada metrica (avg\_ms, p95\_ms, err\_pct, rps, cpu\_mcores, mem\_mib) incluir:
\begin{itemize}
  \item Source
  \item df
  \item sum\_sq
  \item mean\_sq
  \item F
  \item p
  \item eta\_p2
\end{itemize}

\subsection{D.3 Archivos fuente sugeridos}
\begin{itemize}
  \item Testing/results/scaling\_tests/anova\_factorial\_full\_blocked\_by\_day\_20260510.csv
  \item Testing/results/scaling\_tests/anova\_factorial\_matrix\_20260510.csv
\end{itemize}

\section{E. Validacion estadistica y calidad de datos}
\label{anx:validacion}

\subsection{E.1 Verificaciones de consistencia}
\begin{itemize}
  \item Cobertura total de celdas: 48/48.
  \item Replicas por celda: 8.
  \item Revisiones de valores atipicos y corridas fallidas.
\end{itemize}

\subsection{E.2 Indicadores de estabilidad}
\begin{itemize}
  \item Coeficiente de variacion por metrica y por celda.
  \item Trazabilidad de corridas repetidas entre dias.
\end{itemize}

\section{F. Protocolo de reproducibilidad}
\label{anx:protocolo}

\subsection{F.1 Flujo de ejecucion}
\begin{enumerate}
  \item Preparacion del entorno y validacion de dependencias.
  \item Despliegue de la variante a evaluar.
  \item Warmup.
  \item Ejecucion de carga k6.
  \item Recoleccion de metricas de sistema.
  \item Cooldown y limpieza.
  \item Registro de artefactos de salida.
\end{enumerate}

\subsection{F.2 Comandos y scripts}
Incluir comandos exactos y scripts de referencia usados en campana.

\section{G. Artefactos de despliegue y configuracion}
\label{anx:artefactos}

Incluir manifiestos y configuraciones:
\begin{itemize}
  \item YAML de Gateway/Ingress.
  \item Politicas de mTLS.
  \item NetworkPolicy (baseline, basic, strict).
  \item Configuracion de Rate Limiting (baseline, moderate, strict).
\end{itemize}

\section{H. Diccionario de variables}
\label{anx:diccionario}

\begin{itemize}
  \item avg\_ms: latencia promedio por corrida (ms).
  \item p95\_ms: percentil 95 de latencia por corrida (ms).
  \item err\_pct: porcentaje de respuestas con error.
  \item rps: solicitudes por segundo.
  \item cpu\_mcores: uso promedio de CPU (millicores).
  \item mem\_mib: uso promedio de memoria (MiB).
  \item login\_success\_total: total de logins validos durante la corrida.
\end{itemize}

\section{I. Evidencia visual de operacion}
\label{anx:evidencia-visual}

Agregar capturas de paneles de Grafana y, cuando aplique, snapshots de estado del cluster durante corridas representativas:
\begin{itemize}
  \item Paneles de latencia y throughput.
  \item Paneles de CPU y memoria.
  \item Paneles de error rate.
  \item Evidencia de preservacion de usuarios logueados.
\end{itemize}

\section{Nota de uso}

Si el documento principal usa archivos separados por capitulo, este archivo puede integrarse con:
\begin{verbatim}
\chapter{Anexos}
\label{cap:anexos}

Este capitulo consolida los artefactos tecnicos que soportan la trazabilidad, reproducibilidad y auditoria del estudio. Los anexos complementan los capitulos principales y permiten verificar configuracion, ejecucion, procesamiento de datos y resultados estadisticos.

\section{A. Configuracion experimental completa}
\label{anx:config}

\subsection{A.1 Entorno de ejecucion}
\begin{itemize}
  \item Hardware anfitrion: AMD Ryzen 5 3600, 16 GB RAM.
  \item Sistema anfitrion: Windows con WSL2.
  \item VM/WSL2 asignado al experimento: 4 vCPU, 12 GB RAM, 2 GB swap.
  \item Sistema base: Ubuntu 20.04.
\end{itemize}

\subsection{A.2 Stack de herramientas}
\begin{itemize}
  \item Kubernetes: MicroK8s.
  \item Pruebas de carga: k6 v0.50.0.
  \item Observabilidad: Prometheus + Grafana.
  \item Orquestacion de corridas: scripts en Bash/Python del repositorio.
\end{itemize}

\subsection{A.3 Parametros operativos}
\begin{itemize}
  \item Niveles de carga: VUS = 1, 5, 10, 20.
  \item Controles evaluados: C1, C2, C3, C4.
  \item Variantes por control: 3.
  \item Replicas por celda: 8.
  \item Bloqueo temporal: 2 dias de ejecucion.
\end{itemize}

\section{B. Matriz de diseno experimental}
\label{anx:matriz-diseno}

\subsection{B.1 Estructura factorial}
Dise\~no factorial completo: $4 \times 3 \times 4 = 48$ celdas experimentales.

\subsection{B.2 Definicion de factores y niveles}
\begin{itemize}
  \item Factor 1 (Control): C1 API Gateway, C2 mTLS, C3 NetworkPolicy, C4 Rate Limiting.
  \item Factor 2 (Variante): 3 niveles por cada control.
  \item Factor 3 (Carga): VUS 1, 5, 10, 20.
\end{itemize}

\subsection{B.3 Matriz codificada de corridas}
Incluir tabla completa con columnas minimas:
\begin{itemize}
  \item id\_run
  \item day\_block
  \item control
  \item variant
  \item vus
  \item replica
  \item seed (si aplica)
\end{itemize}

\section{C. Tablas completas de resultados}
\label{anx:tablas-resultados}

Incluir resultados agregados y por replica para cada control y variante. Columnas recomendadas:
\begin{itemize}
  \item avg\_ms
  \item p95\_ms
  \item err\_pct
  \item rps
  \item cpu\_mcores
  \item mem\_mib
\end{itemize}

\subsection{C.1 Resultados por control}
\begin{itemize}
  \item C1: API Gateway.
  \item C2: mTLS.
  \item C3: NetworkPolicy.
  \item C4: Rate Limiting.
\end{itemize}

\section{D. Salidas ANOVA completas}
\label{anx:anova}

\subsection{D.1 Especificacion del modelo}
\begin{equation*}
y \sim C(\text{control}) * C(\text{variant\_level}) * C(\text{vus}) + C(\text{day})
\end{equation*}

\subsection{D.2 Tablas ANOVA por metrica}
Para cada metrica (avg\_ms, p95\_ms, err\_pct, rps, cpu\_mcores, mem\_mib) incluir:
\begin{itemize}
  \item Source
  \item df
  \item sum\_sq
  \item mean\_sq
  \item F
  \item p
  \item eta\_p2
\end{itemize}

\subsection{D.3 Archivos fuente sugeridos}
\begin{itemize}
  \item Testing/results/scaling\_tests/anova\_factorial\_full\_blocked\_by\_day\_20260510.csv
  \item Testing/results/scaling\_tests/anova\_factorial\_matrix\_20260510.csv
\end{itemize}

\section{E. Validacion estadistica y calidad de datos}
\label{anx:validacion}

\subsection{E.1 Verificaciones de consistencia}
\begin{itemize}
  \item Cobertura total de celdas: 48/48.
  \item Replicas por celda: 8.
  \item Revisiones de valores atipicos y corridas fallidas.
\end{itemize}

\subsection{E.2 Indicadores de estabilidad}
\begin{itemize}
  \item Coeficiente de variacion por metrica y por celda.
  \item Trazabilidad de corridas repetidas entre dias.
\end{itemize}

\section{F. Protocolo de reproducibilidad}
\label{anx:protocolo}

\subsection{F.1 Flujo de ejecucion}
\begin{enumerate}
  \item Preparacion del entorno y validacion de dependencias.
  \item Despliegue de la variante a evaluar.
  \item Warmup.
  \item Ejecucion de carga k6.
  \item Recoleccion de metricas de sistema.
  \item Cooldown y limpieza.
  \item Registro de artefactos de salida.
\end{enumerate}

\subsection{F.2 Comandos y scripts}
Incluir comandos exactos y scripts de referencia usados en campana.

\section{G. Artefactos de despliegue y configuracion}
\label{anx:artefactos}

Incluir manifiestos y configuraciones:
\begin{itemize}
  \item YAML de Gateway/Ingress.
  \item Politicas de mTLS.
  \item NetworkPolicy (baseline, basic, strict).
  \item Configuracion de Rate Limiting (baseline, moderate, strict).
\end{itemize}

\section{H. Diccionario de variables}
\label{anx:diccionario}

\begin{itemize}
  \item avg\_ms: latencia promedio por corrida (ms).
  \item p95\_ms: percentil 95 de latencia por corrida (ms).
  \item err\_pct: porcentaje de respuestas con error.
  \item rps: solicitudes por segundo.
  \item cpu\_mcores: uso promedio de CPU (millicores).
  \item mem\_mib: uso promedio de memoria (MiB).
  \item login\_success\_total: total de logins validos durante la corrida.
\end{itemize}

\section{I. Evidencia visual de operacion}
\label{anx:evidencia-visual}

Agregar capturas de paneles de Grafana y, cuando aplique, snapshots de estado del cluster durante corridas representativas:
\begin{itemize}
  \item Paneles de latencia y throughput.
  \item Paneles de CPU y memoria.
  \item Paneles de error rate.
  \item Evidencia de preservacion de usuarios logueados.
\end{itemize}

\section{Nota de uso}

Si el documento principal usa archivos separados por capitulo, este archivo puede integrarse con:
\begin{verbatim}
\input{documentecionFinal/anexos}
\end{verbatim}

o bien, si el archivo principal esta en la misma carpeta:
\begin{verbatim}
\input{anexos}
\end{verbatim}

\end{verbatim}

o bien, si el archivo principal esta en la misma carpeta:
\begin{verbatim}
\chapter{Anexos}
\label{cap:anexos}

Este capitulo consolida los artefactos tecnicos que soportan la trazabilidad, reproducibilidad y auditoria del estudio. Los anexos complementan los capitulos principales y permiten verificar configuracion, ejecucion, procesamiento de datos y resultados estadisticos.

\section{A. Configuracion experimental completa}
\label{anx:config}

\subsection{A.1 Entorno de ejecucion}
\begin{itemize}
  \item Hardware anfitrion: AMD Ryzen 5 3600, 16 GB RAM.
  \item Sistema anfitrion: Windows con WSL2.
  \item VM/WSL2 asignado al experimento: 4 vCPU, 12 GB RAM, 2 GB swap.
  \item Sistema base: Ubuntu 20.04.
\end{itemize}

\subsection{A.2 Stack de herramientas}
\begin{itemize}
  \item Kubernetes: MicroK8s.
  \item Pruebas de carga: k6 v0.50.0.
  \item Observabilidad: Prometheus + Grafana.
  \item Orquestacion de corridas: scripts en Bash/Python del repositorio.
\end{itemize}

\subsection{A.3 Parametros operativos}
\begin{itemize}
  \item Niveles de carga: VUS = 1, 5, 10, 20.
  \item Controles evaluados: C1, C2, C3, C4.
  \item Variantes por control: 3.
  \item Replicas por celda: 8.
  \item Bloqueo temporal: 2 dias de ejecucion.
\end{itemize}

\section{B. Matriz de diseno experimental}
\label{anx:matriz-diseno}

\subsection{B.1 Estructura factorial}
Dise\~no factorial completo: $4 \times 3 \times 4 = 48$ celdas experimentales.

\subsection{B.2 Definicion de factores y niveles}
\begin{itemize}
  \item Factor 1 (Control): C1 API Gateway, C2 mTLS, C3 NetworkPolicy, C4 Rate Limiting.
  \item Factor 2 (Variante): 3 niveles por cada control.
  \item Factor 3 (Carga): VUS 1, 5, 10, 20.
\end{itemize}

\subsection{B.3 Matriz codificada de corridas}
Incluir tabla completa con columnas minimas:
\begin{itemize}
  \item id\_run
  \item day\_block
  \item control
  \item variant
  \item vus
  \item replica
  \item seed (si aplica)
\end{itemize}

\section{C. Tablas completas de resultados}
\label{anx:tablas-resultados}

Incluir resultados agregados y por replica para cada control y variante. Columnas recomendadas:
\begin{itemize}
  \item avg\_ms
  \item p95\_ms
  \item err\_pct
  \item rps
  \item cpu\_mcores
  \item mem\_mib
\end{itemize}

\subsection{C.1 Resultados por control}
\begin{itemize}
  \item C1: API Gateway.
  \item C2: mTLS.
  \item C3: NetworkPolicy.
  \item C4: Rate Limiting.
\end{itemize}

\section{D. Salidas ANOVA completas}
\label{anx:anova}

\subsection{D.1 Especificacion del modelo}
\begin{equation*}
y \sim C(\text{control}) * C(\text{variant\_level}) * C(\text{vus}) + C(\text{day})
\end{equation*}

\subsection{D.2 Tablas ANOVA por metrica}
Para cada metrica (avg\_ms, p95\_ms, err\_pct, rps, cpu\_mcores, mem\_mib) incluir:
\begin{itemize}
  \item Source
  \item df
  \item sum\_sq
  \item mean\_sq
  \item F
  \item p
  \item eta\_p2
\end{itemize}

\subsection{D.3 Archivos fuente sugeridos}
\begin{itemize}
  \item Testing/results/scaling\_tests/anova\_factorial\_full\_blocked\_by\_day\_20260510.csv
  \item Testing/results/scaling\_tests/anova\_factorial\_matrix\_20260510.csv
\end{itemize}

\section{E. Validacion estadistica y calidad de datos}
\label{anx:validacion}

\subsection{E.1 Verificaciones de consistencia}
\begin{itemize}
  \item Cobertura total de celdas: 48/48.
  \item Replicas por celda: 8.
  \item Revisiones de valores atipicos y corridas fallidas.
\end{itemize}

\subsection{E.2 Indicadores de estabilidad}
\begin{itemize}
  \item Coeficiente de variacion por metrica y por celda.
  \item Trazabilidad de corridas repetidas entre dias.
\end{itemize}

\section{F. Protocolo de reproducibilidad}
\label{anx:protocolo}

\subsection{F.1 Flujo de ejecucion}
\begin{enumerate}
  \item Preparacion del entorno y validacion de dependencias.
  \item Despliegue de la variante a evaluar.
  \item Warmup.
  \item Ejecucion de carga k6.
  \item Recoleccion de metricas de sistema.
  \item Cooldown y limpieza.
  \item Registro de artefactos de salida.
\end{enumerate}

\subsection{F.2 Comandos y scripts}
Incluir comandos exactos y scripts de referencia usados en campana.

\section{G. Artefactos de despliegue y configuracion}
\label{anx:artefactos}

Incluir manifiestos y configuraciones:
\begin{itemize}
  \item YAML de Gateway/Ingress.
  \item Politicas de mTLS.
  \item NetworkPolicy (baseline, basic, strict).
  \item Configuracion de Rate Limiting (baseline, moderate, strict).
\end{itemize}

\section{H. Diccionario de variables}
\label{anx:diccionario}

\begin{itemize}
  \item avg\_ms: latencia promedio por corrida (ms).
  \item p95\_ms: percentil 95 de latencia por corrida (ms).
  \item err\_pct: porcentaje de respuestas con error.
  \item rps: solicitudes por segundo.
  \item cpu\_mcores: uso promedio de CPU (millicores).
  \item mem\_mib: uso promedio de memoria (MiB).
  \item login\_success\_total: total de logins validos durante la corrida.
\end{itemize}

\section{I. Evidencia visual de operacion}
\label{anx:evidencia-visual}

Agregar capturas de paneles de Grafana y, cuando aplique, snapshots de estado del cluster durante corridas representativas:
\begin{itemize}
  \item Paneles de latencia y throughput.
  \item Paneles de CPU y memoria.
  \item Paneles de error rate.
  \item Evidencia de preservacion de usuarios logueados.
\end{itemize}

\section{Nota de uso}

Si el documento principal usa archivos separados por capitulo, este archivo puede integrarse con:
\begin{verbatim}
\input{documentecionFinal/anexos}
\end{verbatim}

o bien, si el archivo principal esta en la misma carpeta:
\begin{verbatim}
\input{anexos}
\end{verbatim}

\end{verbatim}

\end{verbatim}

\end{verbatim}

o bien, si el archivo principal esta en la misma carpeta:
\begin{verbatim}
\chapter{Anexos}
\label{cap:anexos}

Este capitulo consolida los artefactos tecnicos que soportan la trazabilidad, reproducibilidad y auditoria del estudio. Los anexos complementan los capitulos principales y permiten verificar configuracion, ejecucion, procesamiento de datos y resultados estadisticos.

\section{A. Configuracion experimental completa}
\label{anx:config}

\subsection{A.1 Entorno de ejecucion}
\begin{itemize}
  \item Hardware anfitrion: AMD Ryzen 5 3600, 16 GB RAM.
  \item Sistema anfitrion: Windows con WSL2.
  \item VM/WSL2 asignado al experimento: 4 vCPU, 12 GB RAM, 2 GB swap.
  \item Sistema base: Ubuntu 20.04.
\end{itemize}

\subsection{A.2 Stack de herramientas}
\begin{itemize}
  \item Kubernetes: MicroK8s.
  \item Pruebas de carga: k6 v0.50.0.
  \item Observabilidad: Prometheus + Grafana.
  \item Orquestacion de corridas: scripts en Bash/Python del repositorio.
\end{itemize}

\subsection{A.3 Parametros operativos}
\begin{itemize}
  \item Niveles de carga: VUS = 1, 5, 10, 20.
  \item Controles evaluados: C1, C2, C3, C4.
  \item Variantes por control: 3.
  \item Replicas por celda: 8.
  \item Bloqueo temporal: 2 dias de ejecucion.
\end{itemize}

\section{B. Matriz de diseno experimental}
\label{anx:matriz-diseno}

\subsection{B.1 Estructura factorial}
Dise\~no factorial completo: $4 \times 3 \times 4 = 48$ celdas experimentales.

\subsection{B.2 Definicion de factores y niveles}
\begin{itemize}
  \item Factor 1 (Control): C1 API Gateway, C2 mTLS, C3 NetworkPolicy, C4 Rate Limiting.
  \item Factor 2 (Variante): 3 niveles por cada control.
  \item Factor 3 (Carga): VUS 1, 5, 10, 20.
\end{itemize}

\subsection{B.3 Matriz codificada de corridas}
Incluir tabla completa con columnas minimas:
\begin{itemize}
  \item id\_run
  \item day\_block
  \item control
  \item variant
  \item vus
  \item replica
  \item seed (si aplica)
\end{itemize}

\section{C. Tablas completas de resultados}
\label{anx:tablas-resultados}

Incluir resultados agregados y por replica para cada control y variante. Columnas recomendadas:
\begin{itemize}
  \item avg\_ms
  \item p95\_ms
  \item err\_pct
  \item rps
  \item cpu\_mcores
  \item mem\_mib
\end{itemize}

\subsection{C.1 Resultados por control}
\begin{itemize}
  \item C1: API Gateway.
  \item C2: mTLS.
  \item C3: NetworkPolicy.
  \item C4: Rate Limiting.
\end{itemize}

\section{D. Salidas ANOVA completas}
\label{anx:anova}

\subsection{D.1 Especificacion del modelo}
\begin{equation*}
y \sim C(\text{control}) * C(\text{variant\_level}) * C(\text{vus}) + C(\text{day})
\end{equation*}

\subsection{D.2 Tablas ANOVA por metrica}
Para cada metrica (avg\_ms, p95\_ms, err\_pct, rps, cpu\_mcores, mem\_mib) incluir:
\begin{itemize}
  \item Source
  \item df
  \item sum\_sq
  \item mean\_sq
  \item F
  \item p
  \item eta\_p2
\end{itemize}

\subsection{D.3 Archivos fuente sugeridos}
\begin{itemize}
  \item Testing/results/scaling\_tests/anova\_factorial\_full\_blocked\_by\_day\_20260510.csv
  \item Testing/results/scaling\_tests/anova\_factorial\_matrix\_20260510.csv
\end{itemize}

\section{E. Validacion estadistica y calidad de datos}
\label{anx:validacion}

\subsection{E.1 Verificaciones de consistencia}
\begin{itemize}
  \item Cobertura total de celdas: 48/48.
  \item Replicas por celda: 8.
  \item Revisiones de valores atipicos y corridas fallidas.
\end{itemize}

\subsection{E.2 Indicadores de estabilidad}
\begin{itemize}
  \item Coeficiente de variacion por metrica y por celda.
  \item Trazabilidad de corridas repetidas entre dias.
\end{itemize}

\section{F. Protocolo de reproducibilidad}
\label{anx:protocolo}

\subsection{F.1 Flujo de ejecucion}
\begin{enumerate}
  \item Preparacion del entorno y validacion de dependencias.
  \item Despliegue de la variante a evaluar.
  \item Warmup.
  \item Ejecucion de carga k6.
  \item Recoleccion de metricas de sistema.
  \item Cooldown y limpieza.
  \item Registro de artefactos de salida.
\end{enumerate}

\subsection{F.2 Comandos y scripts}
Incluir comandos exactos y scripts de referencia usados en campana.

\section{G. Artefactos de despliegue y configuracion}
\label{anx:artefactos}

Incluir manifiestos y configuraciones:
\begin{itemize}
  \item YAML de Gateway/Ingress.
  \item Politicas de mTLS.
  \item NetworkPolicy (baseline, basic, strict).
  \item Configuracion de Rate Limiting (baseline, moderate, strict).
\end{itemize}

\section{H. Diccionario de variables}
\label{anx:diccionario}

\begin{itemize}
  \item avg\_ms: latencia promedio por corrida (ms).
  \item p95\_ms: percentil 95 de latencia por corrida (ms).
  \item err\_pct: porcentaje de respuestas con error.
  \item rps: solicitudes por segundo.
  \item cpu\_mcores: uso promedio de CPU (millicores).
  \item mem\_mib: uso promedio de memoria (MiB).
  \item login\_success\_total: total de logins validos durante la corrida.
\end{itemize}

\section{I. Evidencia visual de operacion}
\label{anx:evidencia-visual}

Agregar capturas de paneles de Grafana y, cuando aplique, snapshots de estado del cluster durante corridas representativas:
\begin{itemize}
  \item Paneles de latencia y throughput.
  \item Paneles de CPU y memoria.
  \item Paneles de error rate.
  \item Evidencia de preservacion de usuarios logueados.
\end{itemize}

\section{Nota de uso}

Si el documento principal usa archivos separados por capitulo, este archivo puede integrarse con:
\begin{verbatim}
\chapter{Anexos}
\label{cap:anexos}

Este capitulo consolida los artefactos tecnicos que soportan la trazabilidad, reproducibilidad y auditoria del estudio. Los anexos complementan los capitulos principales y permiten verificar configuracion, ejecucion, procesamiento de datos y resultados estadisticos.

\section{A. Configuracion experimental completa}
\label{anx:config}

\subsection{A.1 Entorno de ejecucion}
\begin{itemize}
  \item Hardware anfitrion: AMD Ryzen 5 3600, 16 GB RAM.
  \item Sistema anfitrion: Windows con WSL2.
  \item VM/WSL2 asignado al experimento: 4 vCPU, 12 GB RAM, 2 GB swap.
  \item Sistema base: Ubuntu 20.04.
\end{itemize}

\subsection{A.2 Stack de herramientas}
\begin{itemize}
  \item Kubernetes: MicroK8s.
  \item Pruebas de carga: k6 v0.50.0.
  \item Observabilidad: Prometheus + Grafana.
  \item Orquestacion de corridas: scripts en Bash/Python del repositorio.
\end{itemize}

\subsection{A.3 Parametros operativos}
\begin{itemize}
  \item Niveles de carga: VUS = 1, 5, 10, 20.
  \item Controles evaluados: C1, C2, C3, C4.
  \item Variantes por control: 3.
  \item Replicas por celda: 8.
  \item Bloqueo temporal: 2 dias de ejecucion.
\end{itemize}

\section{B. Matriz de diseno experimental}
\label{anx:matriz-diseno}

\subsection{B.1 Estructura factorial}
Dise\~no factorial completo: $4 \times 3 \times 4 = 48$ celdas experimentales.

\subsection{B.2 Definicion de factores y niveles}
\begin{itemize}
  \item Factor 1 (Control): C1 API Gateway, C2 mTLS, C3 NetworkPolicy, C4 Rate Limiting.
  \item Factor 2 (Variante): 3 niveles por cada control.
  \item Factor 3 (Carga): VUS 1, 5, 10, 20.
\end{itemize}

\subsection{B.3 Matriz codificada de corridas}
Incluir tabla completa con columnas minimas:
\begin{itemize}
  \item id\_run
  \item day\_block
  \item control
  \item variant
  \item vus
  \item replica
  \item seed (si aplica)
\end{itemize}

\section{C. Tablas completas de resultados}
\label{anx:tablas-resultados}

Incluir resultados agregados y por replica para cada control y variante. Columnas recomendadas:
\begin{itemize}
  \item avg\_ms
  \item p95\_ms
  \item err\_pct
  \item rps
  \item cpu\_mcores
  \item mem\_mib
\end{itemize}

\subsection{C.1 Resultados por control}
\begin{itemize}
  \item C1: API Gateway.
  \item C2: mTLS.
  \item C3: NetworkPolicy.
  \item C4: Rate Limiting.
\end{itemize}

\section{D. Salidas ANOVA completas}
\label{anx:anova}

\subsection{D.1 Especificacion del modelo}
\begin{equation*}
y \sim C(\text{control}) * C(\text{variant\_level}) * C(\text{vus}) + C(\text{day})
\end{equation*}

\subsection{D.2 Tablas ANOVA por metrica}
Para cada metrica (avg\_ms, p95\_ms, err\_pct, rps, cpu\_mcores, mem\_mib) incluir:
\begin{itemize}
  \item Source
  \item df
  \item sum\_sq
  \item mean\_sq
  \item F
  \item p
  \item eta\_p2
\end{itemize}

\subsection{D.3 Archivos fuente sugeridos}
\begin{itemize}
  \item Testing/results/scaling\_tests/anova\_factorial\_full\_blocked\_by\_day\_20260510.csv
  \item Testing/results/scaling\_tests/anova\_factorial\_matrix\_20260510.csv
\end{itemize}

\section{E. Validacion estadistica y calidad de datos}
\label{anx:validacion}

\subsection{E.1 Verificaciones de consistencia}
\begin{itemize}
  \item Cobertura total de celdas: 48/48.
  \item Replicas por celda: 8.
  \item Revisiones de valores atipicos y corridas fallidas.
\end{itemize}

\subsection{E.2 Indicadores de estabilidad}
\begin{itemize}
  \item Coeficiente de variacion por metrica y por celda.
  \item Trazabilidad de corridas repetidas entre dias.
\end{itemize}

\section{F. Protocolo de reproducibilidad}
\label{anx:protocolo}

\subsection{F.1 Flujo de ejecucion}
\begin{enumerate}
  \item Preparacion del entorno y validacion de dependencias.
  \item Despliegue de la variante a evaluar.
  \item Warmup.
  \item Ejecucion de carga k6.
  \item Recoleccion de metricas de sistema.
  \item Cooldown y limpieza.
  \item Registro de artefactos de salida.
\end{enumerate}

\subsection{F.2 Comandos y scripts}
Incluir comandos exactos y scripts de referencia usados en campana.

\section{G. Artefactos de despliegue y configuracion}
\label{anx:artefactos}

Incluir manifiestos y configuraciones:
\begin{itemize}
  \item YAML de Gateway/Ingress.
  \item Politicas de mTLS.
  \item NetworkPolicy (baseline, basic, strict).
  \item Configuracion de Rate Limiting (baseline, moderate, strict).
\end{itemize}

\section{H. Diccionario de variables}
\label{anx:diccionario}

\begin{itemize}
  \item avg\_ms: latencia promedio por corrida (ms).
  \item p95\_ms: percentil 95 de latencia por corrida (ms).
  \item err\_pct: porcentaje de respuestas con error.
  \item rps: solicitudes por segundo.
  \item cpu\_mcores: uso promedio de CPU (millicores).
  \item mem\_mib: uso promedio de memoria (MiB).
  \item login\_success\_total: total de logins validos durante la corrida.
\end{itemize}

\section{I. Evidencia visual de operacion}
\label{anx:evidencia-visual}

Agregar capturas de paneles de Grafana y, cuando aplique, snapshots de estado del cluster durante corridas representativas:
\begin{itemize}
  \item Paneles de latencia y throughput.
  \item Paneles de CPU y memoria.
  \item Paneles de error rate.
  \item Evidencia de preservacion de usuarios logueados.
\end{itemize}

\section{Nota de uso}

Si el documento principal usa archivos separados por capitulo, este archivo puede integrarse con:
\begin{verbatim}
\chapter{Anexos}
\label{cap:anexos}

Este capitulo consolida los artefactos tecnicos que soportan la trazabilidad, reproducibilidad y auditoria del estudio. Los anexos complementan los capitulos principales y permiten verificar configuracion, ejecucion, procesamiento de datos y resultados estadisticos.

\section{A. Configuracion experimental completa}
\label{anx:config}

\subsection{A.1 Entorno de ejecucion}
\begin{itemize}
  \item Hardware anfitrion: AMD Ryzen 5 3600, 16 GB RAM.
  \item Sistema anfitrion: Windows con WSL2.
  \item VM/WSL2 asignado al experimento: 4 vCPU, 12 GB RAM, 2 GB swap.
  \item Sistema base: Ubuntu 20.04.
\end{itemize}

\subsection{A.2 Stack de herramientas}
\begin{itemize}
  \item Kubernetes: MicroK8s.
  \item Pruebas de carga: k6 v0.50.0.
  \item Observabilidad: Prometheus + Grafana.
  \item Orquestacion de corridas: scripts en Bash/Python del repositorio.
\end{itemize}

\subsection{A.3 Parametros operativos}
\begin{itemize}
  \item Niveles de carga: VUS = 1, 5, 10, 20.
  \item Controles evaluados: C1, C2, C3, C4.
  \item Variantes por control: 3.
  \item Replicas por celda: 8.
  \item Bloqueo temporal: 2 dias de ejecucion.
\end{itemize}

\section{B. Matriz de diseno experimental}
\label{anx:matriz-diseno}

\subsection{B.1 Estructura factorial}
Dise\~no factorial completo: $4 \times 3 \times 4 = 48$ celdas experimentales.

\subsection{B.2 Definicion de factores y niveles}
\begin{itemize}
  \item Factor 1 (Control): C1 API Gateway, C2 mTLS, C3 NetworkPolicy, C4 Rate Limiting.
  \item Factor 2 (Variante): 3 niveles por cada control.
  \item Factor 3 (Carga): VUS 1, 5, 10, 20.
\end{itemize}

\subsection{B.3 Matriz codificada de corridas}
Incluir tabla completa con columnas minimas:
\begin{itemize}
  \item id\_run
  \item day\_block
  \item control
  \item variant
  \item vus
  \item replica
  \item seed (si aplica)
\end{itemize}

\section{C. Tablas completas de resultados}
\label{anx:tablas-resultados}

Incluir resultados agregados y por replica para cada control y variante. Columnas recomendadas:
\begin{itemize}
  \item avg\_ms
  \item p95\_ms
  \item err\_pct
  \item rps
  \item cpu\_mcores
  \item mem\_mib
\end{itemize}

\subsection{C.1 Resultados por control}
\begin{itemize}
  \item C1: API Gateway.
  \item C2: mTLS.
  \item C3: NetworkPolicy.
  \item C4: Rate Limiting.
\end{itemize}

\section{D. Salidas ANOVA completas}
\label{anx:anova}

\subsection{D.1 Especificacion del modelo}
\begin{equation*}
y \sim C(\text{control}) * C(\text{variant\_level}) * C(\text{vus}) + C(\text{day})
\end{equation*}

\subsection{D.2 Tablas ANOVA por metrica}
Para cada metrica (avg\_ms, p95\_ms, err\_pct, rps, cpu\_mcores, mem\_mib) incluir:
\begin{itemize}
  \item Source
  \item df
  \item sum\_sq
  \item mean\_sq
  \item F
  \item p
  \item eta\_p2
\end{itemize}

\subsection{D.3 Archivos fuente sugeridos}
\begin{itemize}
  \item Testing/results/scaling\_tests/anova\_factorial\_full\_blocked\_by\_day\_20260510.csv
  \item Testing/results/scaling\_tests/anova\_factorial\_matrix\_20260510.csv
\end{itemize}

\section{E. Validacion estadistica y calidad de datos}
\label{anx:validacion}

\subsection{E.1 Verificaciones de consistencia}
\begin{itemize}
  \item Cobertura total de celdas: 48/48.
  \item Replicas por celda: 8.
  \item Revisiones de valores atipicos y corridas fallidas.
\end{itemize}

\subsection{E.2 Indicadores de estabilidad}
\begin{itemize}
  \item Coeficiente de variacion por metrica y por celda.
  \item Trazabilidad de corridas repetidas entre dias.
\end{itemize}

\section{F. Protocolo de reproducibilidad}
\label{anx:protocolo}

\subsection{F.1 Flujo de ejecucion}
\begin{enumerate}
  \item Preparacion del entorno y validacion de dependencias.
  \item Despliegue de la variante a evaluar.
  \item Warmup.
  \item Ejecucion de carga k6.
  \item Recoleccion de metricas de sistema.
  \item Cooldown y limpieza.
  \item Registro de artefactos de salida.
\end{enumerate}

\subsection{F.2 Comandos y scripts}
Incluir comandos exactos y scripts de referencia usados en campana.

\section{G. Artefactos de despliegue y configuracion}
\label{anx:artefactos}

Incluir manifiestos y configuraciones:
\begin{itemize}
  \item YAML de Gateway/Ingress.
  \item Politicas de mTLS.
  \item NetworkPolicy (baseline, basic, strict).
  \item Configuracion de Rate Limiting (baseline, moderate, strict).
\end{itemize}

\section{H. Diccionario de variables}
\label{anx:diccionario}

\begin{itemize}
  \item avg\_ms: latencia promedio por corrida (ms).
  \item p95\_ms: percentil 95 de latencia por corrida (ms).
  \item err\_pct: porcentaje de respuestas con error.
  \item rps: solicitudes por segundo.
  \item cpu\_mcores: uso promedio de CPU (millicores).
  \item mem\_mib: uso promedio de memoria (MiB).
  \item login\_success\_total: total de logins validos durante la corrida.
\end{itemize}

\section{I. Evidencia visual de operacion}
\label{anx:evidencia-visual}

Agregar capturas de paneles de Grafana y, cuando aplique, snapshots de estado del cluster durante corridas representativas:
\begin{itemize}
  \item Paneles de latencia y throughput.
  \item Paneles de CPU y memoria.
  \item Paneles de error rate.
  \item Evidencia de preservacion de usuarios logueados.
\end{itemize}

\section{Nota de uso}

Si el documento principal usa archivos separados por capitulo, este archivo puede integrarse con:
\begin{verbatim}
\input{documentecionFinal/anexos}
\end{verbatim}

o bien, si el archivo principal esta en la misma carpeta:
\begin{verbatim}
\input{anexos}
\end{verbatim}

\end{verbatim}

o bien, si el archivo principal esta en la misma carpeta:
\begin{verbatim}
\chapter{Anexos}
\label{cap:anexos}

Este capitulo consolida los artefactos tecnicos que soportan la trazabilidad, reproducibilidad y auditoria del estudio. Los anexos complementan los capitulos principales y permiten verificar configuracion, ejecucion, procesamiento de datos y resultados estadisticos.

\section{A. Configuracion experimental completa}
\label{anx:config}

\subsection{A.1 Entorno de ejecucion}
\begin{itemize}
  \item Hardware anfitrion: AMD Ryzen 5 3600, 16 GB RAM.
  \item Sistema anfitrion: Windows con WSL2.
  \item VM/WSL2 asignado al experimento: 4 vCPU, 12 GB RAM, 2 GB swap.
  \item Sistema base: Ubuntu 20.04.
\end{itemize}

\subsection{A.2 Stack de herramientas}
\begin{itemize}
  \item Kubernetes: MicroK8s.
  \item Pruebas de carga: k6 v0.50.0.
  \item Observabilidad: Prometheus + Grafana.
  \item Orquestacion de corridas: scripts en Bash/Python del repositorio.
\end{itemize}

\subsection{A.3 Parametros operativos}
\begin{itemize}
  \item Niveles de carga: VUS = 1, 5, 10, 20.
  \item Controles evaluados: C1, C2, C3, C4.
  \item Variantes por control: 3.
  \item Replicas por celda: 8.
  \item Bloqueo temporal: 2 dias de ejecucion.
\end{itemize}

\section{B. Matriz de diseno experimental}
\label{anx:matriz-diseno}

\subsection{B.1 Estructura factorial}
Dise\~no factorial completo: $4 \times 3 \times 4 = 48$ celdas experimentales.

\subsection{B.2 Definicion de factores y niveles}
\begin{itemize}
  \item Factor 1 (Control): C1 API Gateway, C2 mTLS, C3 NetworkPolicy, C4 Rate Limiting.
  \item Factor 2 (Variante): 3 niveles por cada control.
  \item Factor 3 (Carga): VUS 1, 5, 10, 20.
\end{itemize}

\subsection{B.3 Matriz codificada de corridas}
Incluir tabla completa con columnas minimas:
\begin{itemize}
  \item id\_run
  \item day\_block
  \item control
  \item variant
  \item vus
  \item replica
  \item seed (si aplica)
\end{itemize}

\section{C. Tablas completas de resultados}
\label{anx:tablas-resultados}

Incluir resultados agregados y por replica para cada control y variante. Columnas recomendadas:
\begin{itemize}
  \item avg\_ms
  \item p95\_ms
  \item err\_pct
  \item rps
  \item cpu\_mcores
  \item mem\_mib
\end{itemize}

\subsection{C.1 Resultados por control}
\begin{itemize}
  \item C1: API Gateway.
  \item C2: mTLS.
  \item C3: NetworkPolicy.
  \item C4: Rate Limiting.
\end{itemize}

\section{D. Salidas ANOVA completas}
\label{anx:anova}

\subsection{D.1 Especificacion del modelo}
\begin{equation*}
y \sim C(\text{control}) * C(\text{variant\_level}) * C(\text{vus}) + C(\text{day})
\end{equation*}

\subsection{D.2 Tablas ANOVA por metrica}
Para cada metrica (avg\_ms, p95\_ms, err\_pct, rps, cpu\_mcores, mem\_mib) incluir:
\begin{itemize}
  \item Source
  \item df
  \item sum\_sq
  \item mean\_sq
  \item F
  \item p
  \item eta\_p2
\end{itemize}

\subsection{D.3 Archivos fuente sugeridos}
\begin{itemize}
  \item Testing/results/scaling\_tests/anova\_factorial\_full\_blocked\_by\_day\_20260510.csv
  \item Testing/results/scaling\_tests/anova\_factorial\_matrix\_20260510.csv
\end{itemize}

\section{E. Validacion estadistica y calidad de datos}
\label{anx:validacion}

\subsection{E.1 Verificaciones de consistencia}
\begin{itemize}
  \item Cobertura total de celdas: 48/48.
  \item Replicas por celda: 8.
  \item Revisiones de valores atipicos y corridas fallidas.
\end{itemize}

\subsection{E.2 Indicadores de estabilidad}
\begin{itemize}
  \item Coeficiente de variacion por metrica y por celda.
  \item Trazabilidad de corridas repetidas entre dias.
\end{itemize}

\section{F. Protocolo de reproducibilidad}
\label{anx:protocolo}

\subsection{F.1 Flujo de ejecucion}
\begin{enumerate}
  \item Preparacion del entorno y validacion de dependencias.
  \item Despliegue de la variante a evaluar.
  \item Warmup.
  \item Ejecucion de carga k6.
  \item Recoleccion de metricas de sistema.
  \item Cooldown y limpieza.
  \item Registro de artefactos de salida.
\end{enumerate}

\subsection{F.2 Comandos y scripts}
Incluir comandos exactos y scripts de referencia usados en campana.

\section{G. Artefactos de despliegue y configuracion}
\label{anx:artefactos}

Incluir manifiestos y configuraciones:
\begin{itemize}
  \item YAML de Gateway/Ingress.
  \item Politicas de mTLS.
  \item NetworkPolicy (baseline, basic, strict).
  \item Configuracion de Rate Limiting (baseline, moderate, strict).
\end{itemize}

\section{H. Diccionario de variables}
\label{anx:diccionario}

\begin{itemize}
  \item avg\_ms: latencia promedio por corrida (ms).
  \item p95\_ms: percentil 95 de latencia por corrida (ms).
  \item err\_pct: porcentaje de respuestas con error.
  \item rps: solicitudes por segundo.
  \item cpu\_mcores: uso promedio de CPU (millicores).
  \item mem\_mib: uso promedio de memoria (MiB).
  \item login\_success\_total: total de logins validos durante la corrida.
\end{itemize}

\section{I. Evidencia visual de operacion}
\label{anx:evidencia-visual}

Agregar capturas de paneles de Grafana y, cuando aplique, snapshots de estado del cluster durante corridas representativas:
\begin{itemize}
  \item Paneles de latencia y throughput.
  \item Paneles de CPU y memoria.
  \item Paneles de error rate.
  \item Evidencia de preservacion de usuarios logueados.
\end{itemize}

\section{Nota de uso}

Si el documento principal usa archivos separados por capitulo, este archivo puede integrarse con:
\begin{verbatim}
\input{documentecionFinal/anexos}
\end{verbatim}

o bien, si el archivo principal esta en la misma carpeta:
\begin{verbatim}
\input{anexos}
\end{verbatim}

\end{verbatim}

\end{verbatim}

o bien, si el archivo principal esta en la misma carpeta:
\begin{verbatim}
\chapter{Anexos}
\label{cap:anexos}

Este capitulo consolida los artefactos tecnicos que soportan la trazabilidad, reproducibilidad y auditoria del estudio. Los anexos complementan los capitulos principales y permiten verificar configuracion, ejecucion, procesamiento de datos y resultados estadisticos.

\section{A. Configuracion experimental completa}
\label{anx:config}

\subsection{A.1 Entorno de ejecucion}
\begin{itemize}
  \item Hardware anfitrion: AMD Ryzen 5 3600, 16 GB RAM.
  \item Sistema anfitrion: Windows con WSL2.
  \item VM/WSL2 asignado al experimento: 4 vCPU, 12 GB RAM, 2 GB swap.
  \item Sistema base: Ubuntu 20.04.
\end{itemize}

\subsection{A.2 Stack de herramientas}
\begin{itemize}
  \item Kubernetes: MicroK8s.
  \item Pruebas de carga: k6 v0.50.0.
  \item Observabilidad: Prometheus + Grafana.
  \item Orquestacion de corridas: scripts en Bash/Python del repositorio.
\end{itemize}

\subsection{A.3 Parametros operativos}
\begin{itemize}
  \item Niveles de carga: VUS = 1, 5, 10, 20.
  \item Controles evaluados: C1, C2, C3, C4.
  \item Variantes por control: 3.
  \item Replicas por celda: 8.
  \item Bloqueo temporal: 2 dias de ejecucion.
\end{itemize}

\section{B. Matriz de diseno experimental}
\label{anx:matriz-diseno}

\subsection{B.1 Estructura factorial}
Dise\~no factorial completo: $4 \times 3 \times 4 = 48$ celdas experimentales.

\subsection{B.2 Definicion de factores y niveles}
\begin{itemize}
  \item Factor 1 (Control): C1 API Gateway, C2 mTLS, C3 NetworkPolicy, C4 Rate Limiting.
  \item Factor 2 (Variante): 3 niveles por cada control.
  \item Factor 3 (Carga): VUS 1, 5, 10, 20.
\end{itemize}

\subsection{B.3 Matriz codificada de corridas}
Incluir tabla completa con columnas minimas:
\begin{itemize}
  \item id\_run
  \item day\_block
  \item control
  \item variant
  \item vus
  \item replica
  \item seed (si aplica)
\end{itemize}

\section{C. Tablas completas de resultados}
\label{anx:tablas-resultados}

Incluir resultados agregados y por replica para cada control y variante. Columnas recomendadas:
\begin{itemize}
  \item avg\_ms
  \item p95\_ms
  \item err\_pct
  \item rps
  \item cpu\_mcores
  \item mem\_mib
\end{itemize}

\subsection{C.1 Resultados por control}
\begin{itemize}
  \item C1: API Gateway.
  \item C2: mTLS.
  \item C3: NetworkPolicy.
  \item C4: Rate Limiting.
\end{itemize}

\section{D. Salidas ANOVA completas}
\label{anx:anova}

\subsection{D.1 Especificacion del modelo}
\begin{equation*}
y \sim C(\text{control}) * C(\text{variant\_level}) * C(\text{vus}) + C(\text{day})
\end{equation*}

\subsection{D.2 Tablas ANOVA por metrica}
Para cada metrica (avg\_ms, p95\_ms, err\_pct, rps, cpu\_mcores, mem\_mib) incluir:
\begin{itemize}
  \item Source
  \item df
  \item sum\_sq
  \item mean\_sq
  \item F
  \item p
  \item eta\_p2
\end{itemize}

\subsection{D.3 Archivos fuente sugeridos}
\begin{itemize}
  \item Testing/results/scaling\_tests/anova\_factorial\_full\_blocked\_by\_day\_20260510.csv
  \item Testing/results/scaling\_tests/anova\_factorial\_matrix\_20260510.csv
\end{itemize}

\section{E. Validacion estadistica y calidad de datos}
\label{anx:validacion}

\subsection{E.1 Verificaciones de consistencia}
\begin{itemize}
  \item Cobertura total de celdas: 48/48.
  \item Replicas por celda: 8.
  \item Revisiones de valores atipicos y corridas fallidas.
\end{itemize}

\subsection{E.2 Indicadores de estabilidad}
\begin{itemize}
  \item Coeficiente de variacion por metrica y por celda.
  \item Trazabilidad de corridas repetidas entre dias.
\end{itemize}

\section{F. Protocolo de reproducibilidad}
\label{anx:protocolo}

\subsection{F.1 Flujo de ejecucion}
\begin{enumerate}
  \item Preparacion del entorno y validacion de dependencias.
  \item Despliegue de la variante a evaluar.
  \item Warmup.
  \item Ejecucion de carga k6.
  \item Recoleccion de metricas de sistema.
  \item Cooldown y limpieza.
  \item Registro de artefactos de salida.
\end{enumerate}

\subsection{F.2 Comandos y scripts}
Incluir comandos exactos y scripts de referencia usados en campana.

\section{G. Artefactos de despliegue y configuracion}
\label{anx:artefactos}

Incluir manifiestos y configuraciones:
\begin{itemize}
  \item YAML de Gateway/Ingress.
  \item Politicas de mTLS.
  \item NetworkPolicy (baseline, basic, strict).
  \item Configuracion de Rate Limiting (baseline, moderate, strict).
\end{itemize}

\section{H. Diccionario de variables}
\label{anx:diccionario}

\begin{itemize}
  \item avg\_ms: latencia promedio por corrida (ms).
  \item p95\_ms: percentil 95 de latencia por corrida (ms).
  \item err\_pct: porcentaje de respuestas con error.
  \item rps: solicitudes por segundo.
  \item cpu\_mcores: uso promedio de CPU (millicores).
  \item mem\_mib: uso promedio de memoria (MiB).
  \item login\_success\_total: total de logins validos durante la corrida.
\end{itemize}

\section{I. Evidencia visual de operacion}
\label{anx:evidencia-visual}

Agregar capturas de paneles de Grafana y, cuando aplique, snapshots de estado del cluster durante corridas representativas:
\begin{itemize}
  \item Paneles de latencia y throughput.
  \item Paneles de CPU y memoria.
  \item Paneles de error rate.
  \item Evidencia de preservacion de usuarios logueados.
\end{itemize}

\section{Nota de uso}

Si el documento principal usa archivos separados por capitulo, este archivo puede integrarse con:
\begin{verbatim}
\chapter{Anexos}
\label{cap:anexos}

Este capitulo consolida los artefactos tecnicos que soportan la trazabilidad, reproducibilidad y auditoria del estudio. Los anexos complementan los capitulos principales y permiten verificar configuracion, ejecucion, procesamiento de datos y resultados estadisticos.

\section{A. Configuracion experimental completa}
\label{anx:config}

\subsection{A.1 Entorno de ejecucion}
\begin{itemize}
  \item Hardware anfitrion: AMD Ryzen 5 3600, 16 GB RAM.
  \item Sistema anfitrion: Windows con WSL2.
  \item VM/WSL2 asignado al experimento: 4 vCPU, 12 GB RAM, 2 GB swap.
  \item Sistema base: Ubuntu 20.04.
\end{itemize}

\subsection{A.2 Stack de herramientas}
\begin{itemize}
  \item Kubernetes: MicroK8s.
  \item Pruebas de carga: k6 v0.50.0.
  \item Observabilidad: Prometheus + Grafana.
  \item Orquestacion de corridas: scripts en Bash/Python del repositorio.
\end{itemize}

\subsection{A.3 Parametros operativos}
\begin{itemize}
  \item Niveles de carga: VUS = 1, 5, 10, 20.
  \item Controles evaluados: C1, C2, C3, C4.
  \item Variantes por control: 3.
  \item Replicas por celda: 8.
  \item Bloqueo temporal: 2 dias de ejecucion.
\end{itemize}

\section{B. Matriz de diseno experimental}
\label{anx:matriz-diseno}

\subsection{B.1 Estructura factorial}
Dise\~no factorial completo: $4 \times 3 \times 4 = 48$ celdas experimentales.

\subsection{B.2 Definicion de factores y niveles}
\begin{itemize}
  \item Factor 1 (Control): C1 API Gateway, C2 mTLS, C3 NetworkPolicy, C4 Rate Limiting.
  \item Factor 2 (Variante): 3 niveles por cada control.
  \item Factor 3 (Carga): VUS 1, 5, 10, 20.
\end{itemize}

\subsection{B.3 Matriz codificada de corridas}
Incluir tabla completa con columnas minimas:
\begin{itemize}
  \item id\_run
  \item day\_block
  \item control
  \item variant
  \item vus
  \item replica
  \item seed (si aplica)
\end{itemize}

\section{C. Tablas completas de resultados}
\label{anx:tablas-resultados}

Incluir resultados agregados y por replica para cada control y variante. Columnas recomendadas:
\begin{itemize}
  \item avg\_ms
  \item p95\_ms
  \item err\_pct
  \item rps
  \item cpu\_mcores
  \item mem\_mib
\end{itemize}

\subsection{C.1 Resultados por control}
\begin{itemize}
  \item C1: API Gateway.
  \item C2: mTLS.
  \item C3: NetworkPolicy.
  \item C4: Rate Limiting.
\end{itemize}

\section{D. Salidas ANOVA completas}
\label{anx:anova}

\subsection{D.1 Especificacion del modelo}
\begin{equation*}
y \sim C(\text{control}) * C(\text{variant\_level}) * C(\text{vus}) + C(\text{day})
\end{equation*}

\subsection{D.2 Tablas ANOVA por metrica}
Para cada metrica (avg\_ms, p95\_ms, err\_pct, rps, cpu\_mcores, mem\_mib) incluir:
\begin{itemize}
  \item Source
  \item df
  \item sum\_sq
  \item mean\_sq
  \item F
  \item p
  \item eta\_p2
\end{itemize}

\subsection{D.3 Archivos fuente sugeridos}
\begin{itemize}
  \item Testing/results/scaling\_tests/anova\_factorial\_full\_blocked\_by\_day\_20260510.csv
  \item Testing/results/scaling\_tests/anova\_factorial\_matrix\_20260510.csv
\end{itemize}

\section{E. Validacion estadistica y calidad de datos}
\label{anx:validacion}

\subsection{E.1 Verificaciones de consistencia}
\begin{itemize}
  \item Cobertura total de celdas: 48/48.
  \item Replicas por celda: 8.
  \item Revisiones de valores atipicos y corridas fallidas.
\end{itemize}

\subsection{E.2 Indicadores de estabilidad}
\begin{itemize}
  \item Coeficiente de variacion por metrica y por celda.
  \item Trazabilidad de corridas repetidas entre dias.
\end{itemize}

\section{F. Protocolo de reproducibilidad}
\label{anx:protocolo}

\subsection{F.1 Flujo de ejecucion}
\begin{enumerate}
  \item Preparacion del entorno y validacion de dependencias.
  \item Despliegue de la variante a evaluar.
  \item Warmup.
  \item Ejecucion de carga k6.
  \item Recoleccion de metricas de sistema.
  \item Cooldown y limpieza.
  \item Registro de artefactos de salida.
\end{enumerate}

\subsection{F.2 Comandos y scripts}
Incluir comandos exactos y scripts de referencia usados en campana.

\section{G. Artefactos de despliegue y configuracion}
\label{anx:artefactos}

Incluir manifiestos y configuraciones:
\begin{itemize}
  \item YAML de Gateway/Ingress.
  \item Politicas de mTLS.
  \item NetworkPolicy (baseline, basic, strict).
  \item Configuracion de Rate Limiting (baseline, moderate, strict).
\end{itemize}

\section{H. Diccionario de variables}
\label{anx:diccionario}

\begin{itemize}
  \item avg\_ms: latencia promedio por corrida (ms).
  \item p95\_ms: percentil 95 de latencia por corrida (ms).
  \item err\_pct: porcentaje de respuestas con error.
  \item rps: solicitudes por segundo.
  \item cpu\_mcores: uso promedio de CPU (millicores).
  \item mem\_mib: uso promedio de memoria (MiB).
  \item login\_success\_total: total de logins validos durante la corrida.
\end{itemize}

\section{I. Evidencia visual de operacion}
\label{anx:evidencia-visual}

Agregar capturas de paneles de Grafana y, cuando aplique, snapshots de estado del cluster durante corridas representativas:
\begin{itemize}
  \item Paneles de latencia y throughput.
  \item Paneles de CPU y memoria.
  \item Paneles de error rate.
  \item Evidencia de preservacion de usuarios logueados.
\end{itemize}

\section{Nota de uso}

Si el documento principal usa archivos separados por capitulo, este archivo puede integrarse con:
\begin{verbatim}
\input{documentecionFinal/anexos}
\end{verbatim}

o bien, si el archivo principal esta en la misma carpeta:
\begin{verbatim}
\input{anexos}
\end{verbatim}

\end{verbatim}

o bien, si el archivo principal esta en la misma carpeta:
\begin{verbatim}
\chapter{Anexos}
\label{cap:anexos}

Este capitulo consolida los artefactos tecnicos que soportan la trazabilidad, reproducibilidad y auditoria del estudio. Los anexos complementan los capitulos principales y permiten verificar configuracion, ejecucion, procesamiento de datos y resultados estadisticos.

\section{A. Configuracion experimental completa}
\label{anx:config}

\subsection{A.1 Entorno de ejecucion}
\begin{itemize}
  \item Hardware anfitrion: AMD Ryzen 5 3600, 16 GB RAM.
  \item Sistema anfitrion: Windows con WSL2.
  \item VM/WSL2 asignado al experimento: 4 vCPU, 12 GB RAM, 2 GB swap.
  \item Sistema base: Ubuntu 20.04.
\end{itemize}

\subsection{A.2 Stack de herramientas}
\begin{itemize}
  \item Kubernetes: MicroK8s.
  \item Pruebas de carga: k6 v0.50.0.
  \item Observabilidad: Prometheus + Grafana.
  \item Orquestacion de corridas: scripts en Bash/Python del repositorio.
\end{itemize}

\subsection{A.3 Parametros operativos}
\begin{itemize}
  \item Niveles de carga: VUS = 1, 5, 10, 20.
  \item Controles evaluados: C1, C2, C3, C4.
  \item Variantes por control: 3.
  \item Replicas por celda: 8.
  \item Bloqueo temporal: 2 dias de ejecucion.
\end{itemize}

\section{B. Matriz de diseno experimental}
\label{anx:matriz-diseno}

\subsection{B.1 Estructura factorial}
Dise\~no factorial completo: $4 \times 3 \times 4 = 48$ celdas experimentales.

\subsection{B.2 Definicion de factores y niveles}
\begin{itemize}
  \item Factor 1 (Control): C1 API Gateway, C2 mTLS, C3 NetworkPolicy, C4 Rate Limiting.
  \item Factor 2 (Variante): 3 niveles por cada control.
  \item Factor 3 (Carga): VUS 1, 5, 10, 20.
\end{itemize}

\subsection{B.3 Matriz codificada de corridas}
Incluir tabla completa con columnas minimas:
\begin{itemize}
  \item id\_run
  \item day\_block
  \item control
  \item variant
  \item vus
  \item replica
  \item seed (si aplica)
\end{itemize}

\section{C. Tablas completas de resultados}
\label{anx:tablas-resultados}

Incluir resultados agregados y por replica para cada control y variante. Columnas recomendadas:
\begin{itemize}
  \item avg\_ms
  \item p95\_ms
  \item err\_pct
  \item rps
  \item cpu\_mcores
  \item mem\_mib
\end{itemize}

\subsection{C.1 Resultados por control}
\begin{itemize}
  \item C1: API Gateway.
  \item C2: mTLS.
  \item C3: NetworkPolicy.
  \item C4: Rate Limiting.
\end{itemize}

\section{D. Salidas ANOVA completas}
\label{anx:anova}

\subsection{D.1 Especificacion del modelo}
\begin{equation*}
y \sim C(\text{control}) * C(\text{variant\_level}) * C(\text{vus}) + C(\text{day})
\end{equation*}

\subsection{D.2 Tablas ANOVA por metrica}
Para cada metrica (avg\_ms, p95\_ms, err\_pct, rps, cpu\_mcores, mem\_mib) incluir:
\begin{itemize}
  \item Source
  \item df
  \item sum\_sq
  \item mean\_sq
  \item F
  \item p
  \item eta\_p2
\end{itemize}

\subsection{D.3 Archivos fuente sugeridos}
\begin{itemize}
  \item Testing/results/scaling\_tests/anova\_factorial\_full\_blocked\_by\_day\_20260510.csv
  \item Testing/results/scaling\_tests/anova\_factorial\_matrix\_20260510.csv
\end{itemize}

\section{E. Validacion estadistica y calidad de datos}
\label{anx:validacion}

\subsection{E.1 Verificaciones de consistencia}
\begin{itemize}
  \item Cobertura total de celdas: 48/48.
  \item Replicas por celda: 8.
  \item Revisiones de valores atipicos y corridas fallidas.
\end{itemize}

\subsection{E.2 Indicadores de estabilidad}
\begin{itemize}
  \item Coeficiente de variacion por metrica y por celda.
  \item Trazabilidad de corridas repetidas entre dias.
\end{itemize}

\section{F. Protocolo de reproducibilidad}
\label{anx:protocolo}

\subsection{F.1 Flujo de ejecucion}
\begin{enumerate}
  \item Preparacion del entorno y validacion de dependencias.
  \item Despliegue de la variante a evaluar.
  \item Warmup.
  \item Ejecucion de carga k6.
  \item Recoleccion de metricas de sistema.
  \item Cooldown y limpieza.
  \item Registro de artefactos de salida.
\end{enumerate}

\subsection{F.2 Comandos y scripts}
Incluir comandos exactos y scripts de referencia usados en campana.

\section{G. Artefactos de despliegue y configuracion}
\label{anx:artefactos}

Incluir manifiestos y configuraciones:
\begin{itemize}
  \item YAML de Gateway/Ingress.
  \item Politicas de mTLS.
  \item NetworkPolicy (baseline, basic, strict).
  \item Configuracion de Rate Limiting (baseline, moderate, strict).
\end{itemize}

\section{H. Diccionario de variables}
\label{anx:diccionario}

\begin{itemize}
  \item avg\_ms: latencia promedio por corrida (ms).
  \item p95\_ms: percentil 95 de latencia por corrida (ms).
  \item err\_pct: porcentaje de respuestas con error.
  \item rps: solicitudes por segundo.
  \item cpu\_mcores: uso promedio de CPU (millicores).
  \item mem\_mib: uso promedio de memoria (MiB).
  \item login\_success\_total: total de logins validos durante la corrida.
\end{itemize}

\section{I. Evidencia visual de operacion}
\label{anx:evidencia-visual}

Agregar capturas de paneles de Grafana y, cuando aplique, snapshots de estado del cluster durante corridas representativas:
\begin{itemize}
  \item Paneles de latencia y throughput.
  \item Paneles de CPU y memoria.
  \item Paneles de error rate.
  \item Evidencia de preservacion de usuarios logueados.
\end{itemize}

\section{Nota de uso}

Si el documento principal usa archivos separados por capitulo, este archivo puede integrarse con:
\begin{verbatim}
\input{documentecionFinal/anexos}
\end{verbatim}

o bien, si el archivo principal esta en la misma carpeta:
\begin{verbatim}
\input{anexos}
\end{verbatim}

\end{verbatim}

\end{verbatim}

\end{verbatim}
